\chapter{Term Sheet 与谈判要点}

\section{Term Sheet 的法律地位}

Term Sheet（条款清单）可分为\textbf{有约束力}（如保密、费用、排他）与\textbf{无约束力}（商业核心条款）段落。签署前应明确哪些条款在最终协议中\textbf{不得劣化}，避免签署后被对方律师「技术性推翻」。中英文版本并存时须指定\textbf{优先语言}。

\section{经济条款簇}

投资金额、投前/投后估值、股价与稀释、期权池、清算优先（1x non-participating vs participating）、股息（多为累积与否的象征性条款）。\textbf{反稀释}（full ratchet vs weighted average）在下行轮融资中争议最大，需与创始人心理预期和公司后续融资能力一并评估。

\section{治理与控制簇}

董事会席位、观察员席位、否决权清单（预算、并购、关联交易、高管任免）。过多否决权会瘫痪公司；过少则 LP 对资产保护不足。应在「关键少数」事项上集中否决权，日常运营授权管理层。

\section{流动性与协同退出}

\textbf{领售权（drag-along）}迫使小股东参与整体出售；\textbf{随售权（tag-along）}保护小股东参与买方要约。\textbf{优先认购权}与\textbf{共同出售权}影响后续轮次节奏。IPO 场景下需关注\textbf{注册权}与锁定期。

\section{谈判策略：底线与交换}

事先内部分类：\textbf{必须守住}（如核心信息权、重大资产处置否决）、\textbf{可交换}（如观察员人数）、\textbf{可放弃}（如象征性股息）。用次要让步换取对方在估值或关键治理上的让步。记录谈判历史，避免口头承诺遗失；重大让步须经投资委员会书面确认。
